% ===================================================================
% SKILL: İNTERAKTİF TEST VE CEVAP GİZLEME (OVERLAY)
% ===================================================================
% TANIM:
% Bu şablon, ders sırasında öğrencilere soru sorup, cevabı bir sonraki
% tıkla ekrana getirmek (animasyon/overlay) için kullanılır.
% 
% TEMEL MEKANİZMA:
% \only<n->{...} komutu kullanılır. 
% - n: Slaytın kaçıncı adımında görüneceğini belirtir.
% - ->: O adımdan itibaren görünür kalacağını belirtir.
% 
% ÖRNEK SENARYO:
% 1. Adım: Sadece sorular görünür.
% 2. Adım: İlk sorunun cevabı belirir.
% 3. Adım: İkinci sorunun cevabı belirir...
% ===================================================================

\begin{frame}{Hızlı Test: [Test Konusu Buraya]}
    \justifying{}
    Aşağıdaki soruları cevaplayalım:

    \vspace{1em}
    \begin{enumerate}
        % SORU 1: Cevap 2. tıkta gelir (<2->)
        \item \textbf{Soru Metni 1:} \only<2->{\textcolor{red}{\textbf{Cevap 1}}}
        
        % SORU 2: Cevap 3. tıkta gelir (<3->)
        \item \textbf{Soru Metni 2:} \only<3->{\textcolor{red}{\textbf{Cevap 2}}}
        
        % SORU 3: Cevap 4. tıkta gelir (<4->)
        \item \textbf{Soru Metni 3:} \only<4->{\textcolor{red}{\textbf{Cevap 3}}}
        
        % SORU 4: Cevap 5. tıkta gelir (<5->)
        \item \textbf{Soru Metni 4:} \only<5->{\textcolor{red}{\textbf{Cevap 4}}}
        
        % SORU 5: Cevap 6. tıkta gelir (<6->)
        \item \textbf{Soru Metni 5:} \only<6->{\textcolor{red}{\textbf{Cevap 5}}}
    \end{enumerate}

    \vspace{1.5em}
    \centering
    \textit{\small \textbf{İpucu:} Soruları okuduktan sonra cevabı tahmin etmeleri için öğrencilere süre tanıyın.}
\end{frame}
